% Generated by report_gen.py (ai-ee v1) - do not hand-edit.
\documentclass[11pt,a4paper]{article}
\usepackage[margin=2.2cm]{geometry}
\usepackage{graphicx}
\usepackage{booktabs}
\usepackage{longtable}
\usepackage{pdfpages}
\makeatletter
\@ifundefined{KV@Gin@artifact}{\define@key{Gin}{artifact}[]{}}{}
\makeatother
\usepackage{xcolor}
\usepackage[hidelinks]{hyperref}
\setcounter{tocdepth}{1}
\begin{document}
\sloppy
\begin{center}
{\LARGE\bfseries bb-amp --- Design Document}\\[6pt]
{\large ai-ee v1 pipeline}\\[2pt]
generated 2026-08-16 21:28:38
\end{center}
\tableofcontents

\section{Board and Run Metadata}
{\small
\begin{longtable}{lp{11cm}}
\toprule
\textbf{Field} & \textbf{Value} \\
\midrule
\endhead
board & bb-amp \\
workspace & \texttt{boards/bb-amp} \\
phase & P2 \\
gate status & no gates recorded yet \\
state created & 2026-08-16T17:25:38 \\
state updated & 2026-08-16T21:28:31 \\
\bottomrule
\end{longtable}
}
\section{Overview}
\subsection*{bb-amp brief (owner's words, verbatim)}
learning block-basics: a sensor amplifier. Input a differential 0-20 mV signal on a screw terminal, from an off-board bridge sensor with its own excitation. Output 0-3.3 V single-ended on a screw terminal, DC to 1 kHz. Powered from an external 3.3 V rail on its own screw terminal. Workspace boards/bb-amp.

\subsubsection*{Mode (resolved at P0 by state.py mode)}
\begin{itemize}
\item token: \texttt{learning block-basics:}
\item target: \texttt{block-basics}
\item scope tier: \texttt{block-only}
\item binding level: \texttt{canonical} (geometry is an OUTPUT of the placement)
\item stage under study: none
\end{itemize}
\section{Requirements}
\subsection*{bb-amp - requirements}
Source: \texttt{brief/brief.md} (owner's verbatim brief). Written at P0 by the requirements-analyst. Nothing here selects a part, a topology or a value: the unknowns that would decide those are section 9, and they are for the owner.

\subsubsection*{1. Function}
bb-amp is a sensor-amplifier signal chain on a bare board. It takes a differential 0-20 mV signal from an OFF-BOARD bridge sensor (the sensor carries its own excitation, so this board provides NO excitation circuitry), amplifies it, and presents 0-3.3 V single-ended, referenced to board ground, over DC to 1 kHz. It is powered from an external 3.3 V rail. The nominal transfer is 165 V/V - that is 3.3 V / 20 mV restated, not a topology decision; the exact gain follows from the answers to Q5, Q6 and Q10 and belongs to P2/P4.

\textbf{Build mode} (resolved at P0 by \texttt{state.py mode}; do not re-derive):

\begin{quote}\small\ttfamily\raggedright
\textbar{} dial \textbar{} value \textbar{}\\
\textbar{}---\textbar{}---\textbar{}\\
\textbar{} token \textbar{} `learning block-basics:` \textbar{}\\
\textbar{} mode / target \textbar{} `block-basics` \textbar{}\\
\textbar{} scope tier \textbar{} `block-only` \textbar{}\\
\textbar{} binding level \textbar{} `canonical` - geometry is an OUTPUT of the placement \textbar{}\\
\textbar{} stage under study \textbar{} none \textbar{}
\end{quote}
Scope consequences every later phase and both reviewers should read from here:

\begin{itemize}
\item ON the board: the amplifier block's active part(s), exactly the support
components their datasheets require for correct operation at the stated
operating point (decoupling, the reference the amplifier's own single-supply
operation needs, gain-setting components), plus one input interface, one
output interface and the power entry - the three screw terminals the brief
names.
\item NOT on the board, and never a reviewer finding at \texttt{block-only}: protection
(TVS/ESD, reverse polarity, fusing, OVP/OCP), filtering the datasheet does not
require, indicators, test-points beyond the block's own measurement need,
config straps or DNP options, any second rail or second IC the block does not
need, mechanical and enclosure features beyond mounting the bare board.
\item NOT relaxed by the mode: every electrical spec, every gate, the P2/P3
coverage checks and the research they trigger, DFM, the datasheets' own
requirements, and section 8. A mode relaxes geometry, cost and packaging only.
\end{itemize}
\subsubsection*{2. Interfaces}
Three screw terminals, all named by the brief. They are the block's one input, one output and its power entry - in scope at \texttt{block-only}.

\begin{quote}\small\ttfamily\raggedright
\textbar{} ref \textbar{} interface \textbar{} signal \textbar{} electrical \textbar{}\\
\textbar{}---\textbar{}---\textbar{}---\textbar{}---\textbar{}\\
\textbar{} J1 \textbar{} input \textbar{} differential from the off-board bridge sensor \textbar{} 0-20 mV nominal full scale, DC to 1 kHz; common-mode level UNKNOWN -\textgreater{} Q1; pole count -\textgreater{} Q2 \textbar{}\\
\textbar{} J2 \textbar{} output \textbar{} single-ended, referenced to board GND \textbar{} 0-3.3 V nominal full scale, DC to 1 kHz; achievable span -\textgreater{} Q6; load -\textgreater{} Q8 \textbar{}\\
\textbar{} J3 \textbar{} power \textbar{} external 3.3 V rail + GND \textbar{} see section 3 \textbar{}
\end{quote}
\begin{itemize}
\item \textbf{J1}, recommended 3 poles (IN+, IN-, GND) pending Q2. Two poles only is a
floating input: an amplifier needs a DC return for its input bias currents
and a defined common-mode reference, and without them the output is not
predictable. This is the single riskiest unknown on the board.
\item \textbf{J2}, 2 poles (OUT, GND). The output terminal is also the block's
measurement point, so no separate test points are added.
\item \textbf{J3}, 2 poles (+3V3, GND). One ground node on the board; J1/J2/J3 grounds
are the same net.
\item No connector standard is imposed by the brief. ASSUMED: pluggable or fixed
screw terminal blocks on a 3.5 mm or 5.08 mm pitch accepting roughly 24-16 AWG
wire, top side, silkscreened IN+ / IN- / GND / OUT / +3V3. The exact part is
P3's call.
\item ASSUMED: IN+ above IN- gives a positive output (polarity convention).
\item ASSUMED: no ESD/TVS, no series protection and no anti-alias or band-limiting
filter beyond the amplifier's own response - protection and filtering are
excluded at \texttt{block-only}. Response shape at the top end is Q9.
\item ASSUMED: sensor cable 1 m or shorter, shielded twisted pair, shield landed on
J1's GND pole if Q2's default holds. No cable is supplied with the board.
\item The board does NOT excite the bridge and has no excitation output: the brief
states the sensor brings its own.
\end{itemize}
\subsubsection*{3. Power}
\begin{itemize}
\item Single external rail, 3.3 V, entering on J3. No on-board regulation, no second
rail, no battery, no charging, no power switch (a second rail is excluded at
\texttt{block-only}).
\item Rail tolerance and the current this board may draw: UNKNOWN -\textgreater{} Q10. ASSUMED
for planning: 3.3 V +-5\%.
\item \textbf{GUESS (not a spec):} total board current under 10 mA. A precision
single-supply amplifier plus whatever sets its reference is typically a few
mA; P2/P3 replace this guess with the real datasheet numbers.
\item In scope as datasheet-required support, not as "filtering": local decoupling
at every amplifier supply pin as its datasheet specifies, and the reference
node a single-supply amplifier needs in order to define its output zero.
\item Consequence that drives the design: with only 0 V and 3.3 V available, the
rail bounds BOTH the input common-mode range (Q1) and the output swing (Q6).
Those two questions are the single-supply headroom problem, and they are not
relaxable - they are electrical specs.
\end{itemize}
\subsubsection*{4. Environment}
\begin{itemize}
\item Bench use, indoors, 0-50 C operating (block-only default; taken, not asked).
\item No enclosure. No ingress rating. No vibration, shock or drop requirement.
Non-condensing humidity; no conformal coating.
\item Storage/handling: ordinary lab conditions. No altitude or pollution-degree
requirement follows from a 3.3 V board.
\item Accuracy consequence: drift across 0-50 C, not initial error, is what the
accuracy target in Q7 mostly has to buy.
\end{itemize}
\subsubsection*{5. Size and mounting}
\textbf{No HARD cap. No dimension is stated in the brief, and none is imposed here.} Under binding \texttt{canonical} the board size, aspect and outline are OUTPUTS of the placement, so any number that appeared here would be a PREFERENCE only - RELAXABLE (canonical) - and there is nothing to mark, because nothing is stated.

The canonical flow is mandatory and its order is load-bearing:

1. P5 \texttt{board\_init --outline auto} - generous provisional room. \texttt{board\_init} REFUSES a fixed outline under this binding; guessing a final size here binds placement to a number nobody has earned. 2. P6 place to the canonical layout for this block, gate \texttt{place}. 3. \texttt{board\_edit --outline fit --margin M} - the board becomes what the placement needs. Re-run \texttt{planes\_gen} if it GREW. 4. P7 route.

\begin{itemize}
\item Mounting: none required for bench use. Mechanical is excluded at
\texttt{block-only} (mounting the bare board is all that is in scope), so the absence
of mounting holes is not a reviewer finding; if the earned outline leaves
corner room, plain M3 clearance holes may be added at P5/P6.
\item Height: no limit - no enclosure. Terminal blocks are the tallest parts.
\item Layer count is P2's decision under the "fewest honest layers" default; a
ground reference under a 165 V/V DC signal chain is the reason to expect two.
\end{itemize}
\subsubsection*{6. Quantity and budget}
\begin{itemize}
\item Build quantity 5 (block-only default; taken, not asked).
\item No target unit cost is stated. Cost minimal, within the accuracy answered in
Q7: a 12-bit-usable chain and a 16-bit chain are different part classes and
different money, which is why Q7 is asked rather than assumed cheap.
\item Standard JLC process only; no NRE beyond it; no exotic stackup, no controlled
impedance, no gold fingers.
\end{itemize}
\subsubsection*{7. Assembly}
\begin{itemize}
\item JLC PCBA, single-sided (top) assembly (block-only default; taken, not asked).
\item ASSUMED: the SMT parts are placed by JLC; the three screw terminals are
through-hole in most catalogs and JLC's standard SMT service does not place
them - default is that the owner hand-solders three terminal blocks on the top
side, unless P3 finds suitable SMD terminal blocks. Either resolution is
acceptable; P3 records which.
\item Lead-free process, fab-default finish and soldermask. No conformal coating, no
potting, no cable assemblies, no enclosure.
\end{itemize}
\subsubsection*{8. Compliance and safety flags}
Recorded plainly: \textbf{none of the flagged hazard classes apply to this board.}

\begin{quote}\small\ttfamily\raggedright
\textbar{} flag \textbar{} applies? \textbar{} why \textbar{}\\
\textbar{}---\textbar{}---\textbar{}---\textbar{}\\
\textbar{} mains voltage \textbar{} no \textbar{} powered from an external 3.3 V bench rail (SELV); nothing on the board is mains-referenced \textbar{}\\
\textbar{} battery / charging \textbar{} no \textbar{} no battery, no cell holder, no charge circuit \textbar{}\\
\textbar{} motors / inductive loads \textbar{} no \textbar{} output drives a measurement input, not a load \textbar{}\\
\textbar{} voltage above 30 V \textbar{} no \textbar{} 3.3 V is the highest potential anywhere on the board \textbar{}\\
\textbar{} current above 3 A \textbar{} no \textbar{} expected board current is milliamps (see section 3) \textbar{}\\
\textbar{} RF transmit \textbar{} no \textbar{} no radio, no intentional radiator \textbar{}
\end{quote}
Two conditions would re-open this section, and both hang off the off-board sensor that this board cannot see:

\begin{itemize}
\item If the bridge excitation named in Q1 is above 30 V, J1 becomes a
hazardous-voltage interface and section 8 must be re-answered before P2.
\item If the sensor is mounted on equipment whose chassis sits at mains potential,
the input cable carries that potential onto this board; that is an isolation
requirement, not an amplifier requirement, and it is out of this board's
current scope.
\end{itemize}
No mode grants silence here: if either condition is true, say so in the answers and the pipeline stops to re-scope rather than guessing.

\subsubsection*{9. Open questions}
Batched for one round of answers. Each is closed-form with a recommended default; taking every default is a valid answer.

1. {*}{*}What voltage excites the bridge, and does that excitation share this board's ground?{*}{*} Recommended default: 3.3 V excitation, sharing this board's ground, so both input wires sit near 1.65 V - comfortably inside what a 3.3 V-powered amplifier can accept. Why it matters: with 5 V or 10 V excitation (both common on load cells) the input wires sit near 2.5 V or 5 V; 5 V is OUTSIDE the input range of anything powered from 3.3 V and forces a different front end. Above 30 V it also changes section 8.

2. {*}{*}Does the sensor cable bring a ground wire to this board - i.e. should the input terminal have 3 poles (IN+, IN-, GND) rather than 2?{*}{*} Recommended default: yes, 3 poles, with the cable shield landing on that GND pole. Why it matters: with only two wires the input has no DC path for the amplifier's input bias currents and no defined common-mode level, and the output drifts to a rail. Two poles is only safe if the sensor's excitation return is already tied to this board's ground somewhere else - if so, say where.

3. {*}{*}What is the bridge's nominal resistance (the impedance seen looking back into each input)?{*}{*} Recommended default: 350 ohm, the common load-cell value. Anything from 120 ohm to about 10 k is unremarkable for a high-impedance front end; above 10 k the amplifier's own input current starts to add error and narrows the part choice.

4. {*}{*}Must the output track the excitation (ratiometric), or is a fixed amplifier enough (absolute)?{*}{*} Recommended default: absolute - a fixed gain, with accuracy resting on the excitation being stable. Why it matters: a bridge's output is proportional to its excitation, so a 1\% excitation drift is a 1\% reading error in the absolute case. Ratiometric operation removes that, but needs a sense wire from the off-board excitation supply, which is a second input interface the \texttt{block-only} scope does not carry - answering "ratiometric" is a scope change, not a tweak.

5. \textbf{Beyond the nominal 0-20 mV, what can actually appear at the input?} Recommended default: down to about -1 mV (a real bridge rarely reads exactly zero at rest) and up to about +25 mV under overload, with the amplifier expected to clip gracefully and recover - no added protection parts, which are excluded at this scope. Why it matters: if slightly negative inputs are possible, the output must sit at a small positive pedestal (see Q6) or a negative zero reads as a hard 0 V and cannot be calibrated out.

6. {*}{*}A board on a single 3.3 V rail cannot output exactly 0.000 V or exactly 3.300 V - a real amplifier stops short at each end. Is a guaranteed usable output of roughly 0.05 V to 3.25 V acceptable, with full scale defined inside it?{*}{*} Recommended default: yes, accept the reduced span, and place the 0 mV point at a small positive pedestal (about 0.1 V) so Q5's negative readings stay visible. Why it matters: demanding the true rails requires a supply above 3.3 V or a negative rail - a second rail, which this scope excludes - so this is the one place where the brief's "0-3.3 V" and its "3.3 V rail" genuinely conflict, and the owner picks which one bends.

7. {*}{*}How finely does the reading need to be resolved, and is the system calibrated?{*}{*} Recommended default: usable to 12 bits over the output span (about 0.8 mV at the output, about 5 microvolts referred to the sensor), with zero and span calibrated by whatever reads the output, so initial offset and gain error are trimmed and only noise and 0-50 C drift remain. Why it matters: 8 bits makes almost any amplifier acceptable and cheap; 16 bits changes the part class, the error budget and the cost.

8. \textbf{What will the 0-3.3 V output drive, and over how much cable?} Recommended default: a high-impedance input (bench meter or ADC input, above 100 k) on a lead under 1 m and under about 1 nF of capacitance. Why it matters: a capacitive cable or a load that draws real current changes the output stage and can make an otherwise fine amplifier ring.

9. {*}{*}"DC to 1 kHz" - should the response still be flat (within about 1\%) at 1 kHz, or is 1 kHz where it has already fallen by about 30\% (-3 dB)?{*}{*} Recommended default: flat within about 1\% at 1 kHz, which means designing the amplifier's bandwidth several times higher. Why it matters: it sets how fast the amplifier has to be at a gain of 165, which is a part-selection constraint.

10. {*}{*}How accurate is the external 3.3 V rail, and how much current may this board draw from it?{*}{*} Recommended default: 3.3 V +-5\% and a 10 mA budget. Why it matters: a rail that sags to 3.0 V shrinks the usable output span in Q6, and a tight current budget rules out some precision parts.

\subsubsection*{9a. ANSWERS (owner, 2026-08-16) - these are now REQUIREMENTS}
Every question above is answered. Each is also a \texttt{state.py decision} at P0. Downstream phases design to THIS section, not to section 9's defaults.

\begin{quote}\small\ttfamily\raggedright
\textbar{} Q \textbar{} answer \textbar{}\\
\textbar{}---\textbar{}---\textbar{}\\
\textbar{} 1 \textbar{} Excitation {*}{*}3.3 V, sharing this board's ground{*}{*}. Input common-mode sits at {*}{*}\textasciitilde{}1.65 V{*}{*}. \textbar{}\\
\textbar{} 2 \textbar{} {*}{*}3-pole input terminal{*}{*}: IN+, IN-, GND. Cable shield lands on that GND pole. \textbar{}\\
\textbar{} 3 \textbar{} Bridge {*}{*}350 ohm{*}{*} nominal. \textbar{}\\
\textbar{} 4 \textbar{} {*}{*}Absolute{*}{*} gain (NOT ratiometric). No excitation-sense input. \textbar{}\\
\textbar{} 5 \textbar{} Input range to design for: {*}{*}-1 mV to +25 mV{*}{*}; beyond that the amplifier clips gracefully and recovers. No added protection parts. \textbar{}\\
\textbar{} 6 \textbar{} {*}{*}Reduced span with a pedestal.{*}{*} Guaranteed usable output \textasciitilde{}{*}{*}0.05 V to 3.25 V{*}{*}; 0 mV input maps to a small positive pedestal of about {*}{*}0.1 V{*}{*}. The brief's literal 0.000-3.300 V LOSES to the single 3.3 V rail. \textbar{}\\
\textbar{} 7 \textbar{} {*}{*}12-bit usable{*}{*} over the output span: \textasciitilde{}0.8 mV at the output, \textasciitilde{}{*}{*}5 uV referred to the input{*}{*}. Zero and span calibrated downstream, so 0-50 C drift and noise are the real error budget. \textbar{}\\
\textbar{} 8 \textbar{} Output drives {*}{*}\textgreater{}100 kohm{*}{*}, lead {*}{*}\textless{}1 m{*}{*}, {*}{*}\textless{}1 nF{*}{*}. \textbar{}\\
\textbar{} 9 \textbar{} Response {*}{*}flat within \textasciitilde{}1\% at 1 kHz{*}{*} (1 kHz is inside the flat band, not the corner). \textbar{}\\
\textbar{} 10 \textbar{} Rail {*}{*}3.3 V +-5\%{*}{*}, board current budget {*}{*}10 mA{*}{*}. \textbar{}
\end{quote}
\paragraph*{What those answers make binding}
\begin{itemize}
\item \textbf{Transfer function}: \texttt{Vout = Vped + G {*} Vin\_diff}, with \texttt{Vped \textasciitilde{}= 0.1 V}, and
\texttt{Vin\_diff = 20 mV} (full scale) landing INSIDE the guaranteed usable output
top (\textasciitilde{}3.25 V). The gain follows from those two endpoints and is P2's to fix;
it is close to but NOT equal to the brief's nominal 165 V/V, because the
pedestal and the swing limit both eat span.
\item \textbf{Pedestal source}: the \textasciitilde{}0.1 V reference is a real circuit node. Whatever
drives an instrumentation amplifier's REF input must be LOW impedance (or the
amplifier must have a high-impedance REF by construction) - a bare resistor
divider on a classic 3-op-amp in-amp REF pin degrades CMRR directly, and CMRR
is the whole point of this block. P2 must say which mechanism it is using.
\item \textbf{Bandwidth}: for a single-pole rolloff, 1\% droop at 1 kHz puts the -3 dB
corner near \textbf{7 kHz}. At a gain around 150-165 that is roughly {*}{*}1 MHz of
gain-bandwidth{*}{*} in the gain path. Verify the rolloff-shape assumption; this
is the dominant part-selection constraint and it disqualifies the lowest-power
zero-drift in-amps.
\item \textbf{Error budget} (12-bit usable, calibrated, 0-50 C): offset and gain error
are trimmed out, so what must fit in \textasciitilde{}5 uV RTI is \textbf{offset drift x 50 C},
gain drift, CMRR error against the 1.65 V common mode, and integrated noise
from DC to the amplifier's own bandwidth. Noise density and 1/f corner matter
here; initial offset does not.
\item \textbf{Common mode}: 1.65 V, mid-rail, on a 3.3 V supply - benign, and the reason
a single-supply in-amp front end stays viable. The input range is unipolar
(0-20 mV nominal, -1 mV worst case), so the pedestal, not a mid-rail
reference, sets the output zero.
\end{itemize}
\paragraph*{Delegated authority}
The owner delegated the remaining engineering judgment for this run: decide in-run, optimizing for the LEARNING value of a basic sensor front end. H1-H4 are presented as records rather than blocking holds. The run stops at P9 per the \texttt{block-only} default - ordering stays a separate owner decision, so no irreversible or money-spending step is reached.

\section{Architecture}
\subsection*{Decisions of record (state.json)}
{\small
\begin{longtable}{lp{5.2cm}p{7.2cm}}
\toprule
\textbf{Phase} & \textbf{What} & \textbf{Why} \\
\midrule
\endhead
P0 & Build mode: mode block-basics; scope block-only; binding canonical & token 'learning block-basics:' in the brief; contract in reference/build-modes.md \\
P0 & Q1 excitation: 3.3 V, sharing this board's ground -\textgreater{} input common-mode \textasciitilde{}1.65 V & owner answer; keeps the inputs inside what a 3.3 V-powered in-amp accepts, so the canonical single-IC front end stays available (5 V or 10 V excitation would force a difference-amp with input attenuation) \\
P0 & Q6 output span: usable \textasciitilde{}0.05-3.25 V with a small positive pedestal (\textasciitilde{}0.1 V) at 0 mV input; the brief's literal 0.000-3.300 V loses & owner answer; a single 3.3 V rail cannot reach either rail exactly, and the alternative (second rail / charge pump) is excluded at block-only. The pedestal also keeps slightly-negative bridge readings (Q5, -1 mV) visible and calibratable. This is an ELECTRICAL spec resolution, not a mode relaxation \\
P0 & Q7 accuracy: 12-bit usable over the output span (\textasciitilde{}0.8 mV out, \textasciitilde{}5 uV referred to input), zero and span calibrated downstream & owner answer; makes 0-50 C drift and noise the real error budget rather than initial offset, and keeps the part in the mainstream precision in-amp class rather than the 16-bit zero-drift class \\
P0 & Q9 bandwidth: response flat within \textasciitilde{}1\% at 1 kHz (not -3 dB at 1 kHz) & owner answer; for a single-pole rolloff 1\% droop at 1 kHz implies a -3 dB corner near 7 kHz, so at the full-scale gain the amplifier needs roughly 1 MHz gain-bandwidth - this is now the dominant part-selection constraint \\
P0 & Q2/Q3/Q4/Q5/Q8/Q10 take their documented defaults: 3-pole input (IN+/IN-/GND, shield to GND), 350 ohm bridge, absolute (not ratiometric) gain, input -1 mV to +25 mV clipping gracefully, output into \textgreater{}100 k / \textless{}1 m / \textless{}1 nF, rail 3.3 V +-5\% with a 10 mA budget & each default is the low-risk answer and none changes the topology; ratiometric in particular would add an excitation-sense input, i.e. a second input interface, which block-only scope excludes \\
P0 & Owner delegated the remaining engineering judgment: decide in-run, optimizing for the best LEARNING value of a basic sensor front end. H1-H4 are presented as records and the run proceeds without blocking; the run stops at P9 (block-only default) so no irreversible or money-spending step is reached & owner instruction 2026-08-16 ('Make decisions yourself based on what will give the best learning for the basic system') plus three boards running in parallel; nothing before P10 is irreversible, and ordering stays a separate owner decision \\
P1 & P1 roster: research-component-scout (block 'inamp') + research-reference-design (block 'bridge-front-end') only & single functional block. research-power-architect skipped: trivially powered (one external 3.3 V rail, no conversion, no second rail at block-only). research-interface-spec skipped: screw terminals are not standards-bound, so there is no spec to extract \\
P1 & Reconciliation handed to P2: the scout's 'AD8226 REF grounded, pedestal injected at the second stage' does NOT survive the part's guaranteed output floor of (V-)+0.1 V - at G1=100 every input below +2 mV is in first-stage saturation, which is the bottom 10 percent of full scale. A positive, low-impedance-driven REF is unavoidable for any classic in-amp here & the pedestal exists to keep the -1 mV worst-case input visible (requirements 9a Q5/Q6); a stage-2 pedestal cannot recover information stage 1 has already clipped. Once a REF buffer is required anyway, single-stage at G\textasciitilde{}147 uses the same two ICs with fewer passives than the split \\
P2 & Topology: two-stage. U1 AD8226 at G1=39.9 (RG 1.27k, REF on a buffered 0.252 V node) -\textgreater{} U2B OPA2333-class RRIO at G2=3.49 whose gain resistor RETURNS TO /VREF, so Vout = Vref + G1{*}G2{*}Vdiff. U2A buffers the R2/R3 pedestal divider. Total G=139.2 & forced by the AD8226 input-voltage-range inequalities (Rev.C p.20 Eq.1-3 + Table 8, validated against published Fig.12/13 before use): at Vcm=1.65 V on a 3.3 V rail the amplified differential swing at the internal gain-stage nodes is capped at 1.90 V (25 C) and 1.33 V in the rail-5pct/excitation+5pct/0 C corner, so max single-stage gain is 66-95, not 147. The datasheet prescribes exactly this remedy. Returning the stage-2 gain resistor to /VREF makes the output pedestal EXACTLY Vref with no matched network and cancels reference drift to first order \\
P2 & Electrical specs derived away from the brief (NOT mode relaxations): gain 165 -\textgreater{} 139.2; pedestal \textasciitilde{}0.1 V -\textgreater{} 0.252 V; usable output window 0.05/3.25 V -\textgreater{} 0.10/3.20 V; transfer Vout = 0.252 + 139.2{*}Vdiff (-1 mV -\textgreater{} 0.113 V, 20 mV -\textgreater{} 3.037 V, +25 mV clips gracefully) & gain must fit under the -5pct rail ceiling of 3.065 V; the -1 mV worst case of Q5 needs 0.139 V of headroom below zero scale plus the stage output floor. A mode relaxes geometry/cost/packaging only, so each of these is recorded as an electrical derivation with its mechanism \\
P2 & ACCURACY SPEC MISSED AND ACCEPTED: Q7's 5 uV RTI over 0-50 C is not met. Offset drift 13.9 uV typ / 56.4 uV A-grade max over +-25 C from a 25 C calibration, plus 30-87 uV of gain drift at full scale. Pedestal additionally tracks the rail at 548 uV RTI per volt & 5 uV RTI on a 20 mV full scale over 25 C is 10 ppm/degC of full scale - a laboratory-grade tempco, not a 12-bit one - and no 3.3 V-capable part in the candidate set reaches it. The alternative (INA333 at G\textasciitilde{}40 plus a zero-drift second stage) spends the entire bandwidth margin (8.75 kHz vs 45 kHz), lands ON the budget typical-only, and leaves the largest term (gain drift) untouched. Accept-and-record chosen per the owner's learning-value instruction; the miss is quantified here and shown at H1, never dropped \\
P2 & Stackup JLC2313\_1.6, 2 layers, class 2L & no controlled impedance exists on a DC-1 kHz chain; the input pair needs a CONTINUOUS ground reference (solid B.Cu pour, no slot under the differential run), not a close one. 4L would duplicate that job and distribute 0.65 mA \\
P2 & No output RC filter. Chain noise is 0.88 uVrms RTI band-limited to 1 kHz but 4.97 uVrms (33 uV p-p) over the chain's own 41 kHz & an output filter is filtering the datasheets do not require, which block-only scope excludes. Recorded so the number is visible: the reader must band-limit to get the low figure. This is a scope decision, never a claim that the filter is unnecessary in general \\
P2 & AD8227 (pin-compatible, named by the AD8226 datasheet as the other escape from the Eq.2 boundary) deferred to the P2 research task as a directed sub-question rather than settled by assumption & LCSC has it but thin (AD8227ARMZ-R7 MSOP-8, 201 stock, USD 2.99; the SOIC-8 A-grade has 18). Its boundary constants and its bandwidth at G\textasciitilde{}140 are unverified, and the architect correctly refused to design against unverified numbers. The researcher has allowlisted analog.com access, so a cited answer costs nothing extra and becomes a knowledge record either way \\
P2 & Mode relaxed NOTHING on this board & no dimension is stated anywhere in the brief, geometry is an output of the placement under the canonical binding, and cost/packaging were untouched \\
P2 & analog.com is unreachable from this host - whole-host timeout, not just a PDF path (research.py fetch ReadTimeout twice; bare curl to https://www.analog.com/ returns 0 bytes after 90 s). Treat as a STANDING condition, not a blip: P1's reference-design agent hit the same thing independently & consequence for this board: no vendor-layout-tier source exists for the actual part (AD8226), so the interface records are all cross-vendor tier witnessed by TI's INA333 instead of ADI's own text. WORKAROUND that works and should be used from here: the AD8226 datasheet is obtainable through the allowlisted LCSC mirror (wmsc.lcsc.com), and the block-inamp-1 task already has 2304140030\_Analog-Devices-AD8226ARZ-R7\_C34250.pdf in the shared sources dir \\
P2 & The interface:in slot will be RETYPED before the post-research coverage re-run: its operating\_point currently carries only diff-pair geometry (gap\_mm/max\_skew\_mm/max\_uncoupled\_mm/term\_pair\_mm), so none of the five records' envelopes can be tested and they can only ever read provisional & measured by the researcher, not assumed: it ran envelope containment against the slot and all five topology records returned unknown. A kilohertz-rate 20 mV differential sensor input is not a controlled-impedance pair - what it needs declared is rsource\_ohm, vcm\_v, vout\_v, vsig\_mv, dt\_c, err\_rti\_mv, ibias\_ua, mismatch\_pct. diff\_pairs{[}{]}.operating\_point is supported by knowledgelib, so the retype is legal without a schema change. Deferred until the block task closes so the AD8226 input bias current comes from a cited page rather than a guess \\
P2 & AD8227 stays UNCLOSED and the design does NOT switch to it. The two-stage AD8226 chain stands & its datasheet is unacquirable from this host today: analog.com and mouser.com stall research.py fetch, and LCSC's mirror serves a \textasciitilde{}7 KB 'Datasheet temporarily unavailable' stub for every ADI part. Neither its input-range constants at Vs=3.3 V / Vcm=1.65 V nor its bandwidth at G\textasciitilde{}140 were ever read. The derived mechanism points to yes (an in-amp's minimum gain IS its difference-amp gain, so min-gain-5 means one fifth the internal-node swing at equal total gain) but a derivation is not a datasheet, and trading a validated design for an unverifiable part is the wrong trade. Recorded as the highest-value next acquisition, together with ADI AN-1401 'Instrumentation Amplifier Common-Mode Range: The Diamond Plot' \\
P2 & P6/P7 CORRECTION: do NOT spend layout effort on input trace/via length matching and plane-continuity fetishism for CMRR. The rule CMRR=(G+1)/(4t) belongs to a four-discrete-resistor difference amplifier where G and t are the SAME network; an in-amp's \textasciitilde{}100 Gohm inputs convert essentially nothing from series-path mismatch at DC-1 kHz with a 350 ohm source & second reader refutation, page-level: every SBOA582 circuit is a four-resistor diff-amp at G=1, and the transfer to stray trace asymmetry was asserted rather than argued. The error term that IS verified at the 5 uV RTI budget is the TEMPERATURE GRADIENT across the two input junctions - so the placement rule that matters is keeping the two input poles at the same temperature and away from anything that dissipates, not matching their lengths \\
P2 & BRING-UP NOTE: a missing input bias-current return path does NOT park the output at a rail - it produces a plausible near-zero output that 'appears normal' & second reader found the researcher's claim contradicted by INA333 s8.2.2.6 on the same cited page and by TI SLYT226 p.1. This matters because the failure mode is silent: a board with a floating input reads like a working board with no signal. Belongs in the bring-up notes \\
P2 & interface-in-1 ledger quality: 1 of 6 records verified. The refutations are themselves the finding - the whole ledger is TI/Microchip-witnessed and NO number for this board's actual part (AD8226) is witnessed anywhere in it, so the bias-return REQUIREMENT transfers but the SIZING does not & analog.com is unreachable (standing condition), so the interface researcher had no vendor witness for the board's own part. The block-inamp-1 task has since proved the AD8226 datasheet IS obtainable from an allowlisted Farnell mirror - that source did not exist when this task ran \\
P2 & The architecture's foundation is VERIFIED on the page and the generalization is stronger than the architect claimed: AD8226 printed p.19 gives Vout = G(Vin+ - Vin-) + Vref exactly, so Eq.2 becomes Vcm + \textbar{}Vout - Vref\textbar{}/2 \textless{} Vs - V+LIMIT with NO GAIN TERM IN IT. Confirmed on the printed figures, not just algebraically: Fig.9 (G=1) and Fig.12 (G=100) at Vs=+2.7 V differ by at most 0.1 V at four vertices, and the Vref=0 diagonal from (+0.02,+2.0) to (+2.4,+0.8) is the SAME straight line in both - that line IS Eq.2. All twelve Table 8 constants exact & the two-stage topology is downstream of this claim, so it was targeted for refutation by a fresh reader and survived. The practical restatement: the in-amp's constraint is on ITS OWN OUTPUT SWING, not on its gain - which is exactly why moving gain to the second stage relieves it \\
P2 & AMENDING the earlier P6/P7 symmetry ruling. Do NOT chase input trace LENGTH matching for DC CMRR (that stands - an in-amp's \textasciitilde{}100 Gohm inputs convert nothing from series mismatch at DC). But DO keep the two input paths electrically similar - same width, same neighbourhood, same environment - because AC CMRR does degrade with SOURCE IMBALANCE & the block second reader refuted the parasitics record precisely here: its envelope claimed no vendor number exists above 5 kHz, but the ledger's own AD8226 PDF p.13 carries Fig.27 'CMRR vs Frequency, RTI' and Fig.28 'CMRR vs Frequency, RTI, 1 kohm Source Imbalance' out to 100 kHz, and Fig.28 measures this exact mechanism. Our chain is open to 41 kHz and common-mode interference rides in there. Our own imbalance is far below Fig.28's 1 kohm case, so this is cheap insurance, not a driver - the dominant term at the 5 uV budget remains the thermal gradient across the two input junctions \\
P2 & AD8226 A-grade input bias current, cited from Table 3 printed p.5: 20 nA typ / 27 nA max at +25 C (5 nA min), 30/35 nA at -40 C, 15/25 nA at +125 C, avg TC 70 pA/degC. Input offset current 1 nA max. The input stage is pnp, so bias current always flows OUT of the part - the return path must SINK it & needed to declare the input interface's operating point without guessing; obtained by asking the second reader while it had Table 3 open rather than by inventing a number. The pnp direction is a bring-up-relevant detail: the return network sinks current, it does not source it \\
P2 & P6/P7 LAYOUT RULE, now witnessed by the board's own part (AD8226 Fig.28): SOURCE-PATH IMBALANCE, not gain, sets CMRR here. Balanced (Fig.27) LF CMRR climbs 96/117/133/144 dB across G=1/10/100/1000; with a 1 kohm source imbalance it climbs only 96/102/106/108 dB, and above \textasciitilde{}1 kHz the G=1/10/100 curves collapse onto one roll-off reaching 55-68 dB at 100 kHz. At this board's G1=39.9 that imbalance costs about 25 dB at the 1 kHz signal frequency & supersedes the vaguer earlier ruling with a measured number from the actual part. Consequence for this board: it adds NO series component in either input leg, and a 350 ohm bridge's arm-to-arm imbalance is orders below Fig.28's 1 kohm case, so we are far from the penalty - but the rule to enforce at P6/P7 is that nothing may be inserted in one input leg and not the other, and the two paths must see the same environment. That is a different instruction from length-matching, which still buys nothing at DC \\
P2 & Recorded as a research-quality practice worth keeping: the researcher caught that AD8226 Figures 27/28 are taken at Vs=+-15 V (TPC conditions header, printed p.9) while Table 3 is at +Vs=2.7 V single supply, so the absolute dB from those curves cannot be set against Table 3 minimums - only the PENALTY (the ratio of two curves under identical conditions) transfers & without that check the record would have carried a false claim that source imbalance drops the part below its guaranteed CMRR spec. Both records state the limitation explicitly. This is the failure mode adversarial review exists to catch, caught by the author instead \\
P2 & CORRECTION to the P6/P7 CMRR layout rule: the 1 kHz source-imbalance penalty is about 16-17 dB (17.0/15.9/16.3 dB at G=10/100/1000), NOT the \textasciitilde{}25 dB I recorded earlier. The qualitative rule is unchanged - source-path imbalance, not gain, sets CMRR here - only the magnitude was overstated & second-reader refutation with page-level measurements: at 1 kHz Fig.27 reads 109.4/109.4/110.6 dB and Fig.28 reads 92.4/93.5/94.3 dB. The researcher's claimed 27/23/19 dB was 3-10 dB out, outside its own stated +-2 dB curve-read tolerance, and that is exactly the number the rule hands a designer at this board's 1 kHz signal frequency. The reader also verified everything else in the record (plateau ladders within 1 dB, gain-to-curve mapping via the Fig.28 arrow callouts, the +-15 V TPC conditions caveat, envelope bounds) and found a bonus fact: Table 2 (printed p.3) is the matching +-15 V table carrying the SAME minimums as Table 3, so the conditions caveat costs nothing \\
P2 & CORRECTION to the earlier bring-up note. A missing bias-current return path has NO diagnostic symptom: the AD8226 page states none; Renesas AN1298 p.14 says the output will slowly drift until it saturates into a rail; the same vendor's sensor-health paper says a saturated input stage can appear normal to the processing circuitry. Both endings are real. The honest bring-up rule is: PROBE THE INPUT PINS, NOT THE OUTPUT & my earlier note recorded only one of the two endings (plausible near-zero output), inheriting it from round one's TI-witnessed record. The re-witnessing pass found the contradiction is between two real documented behaviours, not an error in either, so the diagnostic value is in the input pins' DC level \\
P2 & Bias-return sizing now witnessed by the board's own part, and it confirms this board needs NO bias resistors: the 350 ohm bridge is excited from a 3.3 V supply sharing our ground, so each input already has a DC return through its own bridge arms and the J1 GND pole. The sizing rule matters only if the source were floating - and it is brutal there: AD8226's own printed value is 10 Mohm (single resistor, Fig.62), but with a single resistor the far input's bias current returns through the 350 ohm source for about 10 uV differential, already past the 5 uV budget. A balanced pair caps R near 3.7 kohm (differential term = Ios{*}R + Ib{*}R{*}mismatch), which then loads a 350 ohm source by about 4.5 pct. Round one's transplanted 47k figure would have cost about 60 uV differential on this part & this is exactly the claim the first reader said would not transfer - the REQUIREMENT transfers, the SIZING does not - and re-witnessing proved it right: the CMOS 47k/10k values were off by more than an order of magnitude for a 27 nA pnp-input part \\
P2 & P6/P7 input-path rule sharpened to a number: with 0.8 Gohm\textbar{}\textbar{}2 pF inputs the DC resistance term is dead, and it is CAPACITANCE match that matters - 1 pF of path mismatch gives -119 dB at 1 kHz, 10 pF gives -85 dB at 5 kHz against the part's own 90 dB spec & gives placement and routing a checkable target (keep the two input paths' capacitance to the pour matched to about 1 pF) instead of the untestable instruction to match lengths. Consistent with, and more actionable than, the Fig.28 source-imbalance rule \\
P2 & P2 exit coverage passes: exit 0, 2 slots, 0 gaps, all 11 classes provisional. No 'designing under coverage gap' decision is needed & 16 verified records and 2 new checklists (inamp, in) now cover both slots. Provisional is the correct ceiling here and not a shortfall: research narrows a gap to provisional by contract, the owner's promotion ruling makes a record approved, and only bring-up evidence makes it proven \\
P2 & Research spend: 4 tasks over 2 rounds, all 6 of the per-run budget. Round 1 (13 records) was witnessed by the wrong parts because analog.com was unreachable - 7 refuted. Round 2 (9 records) re-witnessed from the AD8226 itself via an allowlisted Farnell mirror - 9 verified, 9 of 9 & the second round was worth its cost and the evidence says so: it closed all five interface gap classes, replaced a transplanted 47k bias-return figure that would have cost \textasciitilde{}60 uV differential on this part, and produced a corrected CMRR-vs-imbalance number that a layout rule now rests on. Deciding to spend it rather than record 'designing under coverage gap' was the right call for a board whose stated purpose is to seed the library \\
\bottomrule
\end{longtable}
}
\subsection*{Sheet plan}
\subsection*{bb-amp - schematic sheet plan (P2)}
\subsubsection*{Ruling: ONE sheet. No hierarchy.}
The whole board is 14 components on one signal path (J1 -\textgreater{} U1 -\textgreater{} U2B -\textgreater{} J2) plus a two-resistor reference and four capacitors. A hierarchical sheet buys navigation at the cost of hierarchical pins, sheet instances and a netlist whose names carry sheet paths - all of it overhead against a schematic that fits on one A4 page with room for the design equations in a text box. Splitting it would also hide the one thing a reader of this board should see at a glance: that \texttt{/VREF} is a single node touching U1 pin 6, U2A output and the stage-2 gain return, which is what makes the pedestal exact (\texttt{blocks.md} section 2).

Root sheet = the only sheet. Net names carry a leading \texttt{/} for local labels (\texttt{/IN\_P}), power nets are bare (\texttt{+3V3}, \texttt{GND}).

\subsubsection*{Sheet: root (\texttt{bb-amp.kicad\_sch})}
\textbf{Blocks on it:} B1 input interface, B2 in-amp, B3 reference/pedestal, B4 output gain stage, B5 output interface, B6 power entry.

\textbf{Layout of the drawing} (left to right, matching the board's own chain, so placement groups at P6 read straight off the schematic): J1 / input pair at the far left; U1 with R1 across its RG pins; the reference cluster (R2, R3, C4, U2A) below U1; U2B with R4/R5 to the right of U1; J2 at the far right; J3 and C2 along the bottom with C1/C3 drawn at their IC supply pins.

\textbf{Nets that leave a block} (these become the placement groups at P6 and the net classes at P5):

\begin{quote}\small\ttfamily\raggedright
\textbar{} net \textbar{} from \textbar{} to \textbar{} note \textbar{}\\
\textbar{}---\textbar{}---\textbar{}---\textbar{}---\textbar{}\\
\textbar{} `/IN\_P` \textbar{} J1 pin 1 \textbar{} U1 pin 4 (+IN) \textbar{} differential pair with `/IN\_N`; matched, symmetric, over unbroken B.Cu \textbar{}\\
\textbar{} `/IN\_N` \textbar{} J1 pin 2 \textbar{} U1 pin 1 (-IN) \textbar{} same \textbar{}\\
\textbar{} `GND` \textbar{} J1 pin 3, J2 pin 2, J3 pin 2 \textbar{} plane \textbar{} bias-current return path for both inputs (refdesign D12) - not a convenience pole \textbar{}\\
\textbar{} `/VREF\_SET` \textbar{} R2/R3 junction, C4 \textbar{} U2A + input \textbar{} high-Z divider node, bypassed by C4 \textbar{}\\
\textbar{} `/VREF` \textbar{} U2A output \textbar{} U1 pin 6 (REF), R5 \textbar{} must stay below 2 ohm of source impedance (AD8226 Reference Terminal); one node, three loads \textbar{}\\
\textbar{} `/AMP1\_OUT` \textbar{} U1 pin 7 (OUT) \textbar{} U2B + input \textbar{} 0.21 - 1.05 V over the input range \textbar{}\\
\textbar{} `/FB2` \textbar{} R4/R5 junction \textbar{} U2B - input \textbar{} stage-2 summing node \textbar{}\\
\textbar{} `/VOUT` \textbar{} U2B output \textbar{} J2 pin 1 \textbar{} 0.113 - 3.037 V into \textgreater{}= 100 k \textbar{}\\
\textbar{} `+3V3` \textbar{} J3 pin 1 \textbar{} U1 pin 8, U2 pin 8, R2, C1, C2, C3 \textbar{} 0.65 mA worst case \textbar{}
\end{quote}
\subsubsection*{Refdes ranges}
One sheet, so uniqueness is trivial - the ranges are still fixed here so P4 and any later \texttt{schem\_refdes} renumber land in the same places, and so a second sheet added later cannot collide.

\begin{quote}\small\ttfamily\raggedright
\textbar{} class \textbar{} range \textbar{} assigned on this board \textbar{}\\
\textbar{}---\textbar{}---\textbar{}---\textbar{}\\
\textbar{} U (ICs) \textbar{} U1 - U9 \textbar{} U1 AD8226 in-amp; U2 OPA2333 dual (U2A buffer, U2B gain) \textbar{}\\
\textbar{} R \textbar{} R1 - R19 \textbar{} R1 RG 1.27 k; R2 121 k, R3 10.0 k divider; R4 24.9 k fb, R5 10.0 k gain return \textbar{}\\
\textbar{} C \textbar{} C1 - C19 \textbar{} C1 100 n at U1; C2 10 u bulk; C3 100 n at U2; C4 100 n on `/VREF\_SET` \textbar{}\\
\textbar{} J \textbar{} J1 - J9 \textbar{} J1 input 3-pole; J2 output 2-pole; J3 power 2-pole \textbar{}\\
\textbar{} H (mounting) \textbar{} H1 - H4 \textbar{} none required; optional M3 clearance holes only if the earned outline leaves corner room \textbar{}\\
\textbar{} `\#PWR` \textbar{} `pwr\_base = 100` -\textgreater{} `\#PWR0101...` \textbar{} power/GND symbols on the root sheet \textbar{}\\
\textbar{} `\#FLG` \textbar{} `\#FLG0101...` \textbar{} one PWR\_FLAG on `+3V3` and one on `GND` (external rail, no on-board source - ERC needs them) \textbar{}
\end{quote}
\subsubsection*{Values and tolerances the schematic must carry}
Calibration removes initial error but not tempco, so tolerance and TCR are split deliberately (\texttt{blocks.md} Ruling 3 quantifies each term):

\begin{quote}\small\ttfamily\raggedright
\textbar{} ref \textbar{} value \textbar{} tolerance / TCR \textbar{} why \textbar{}\\
\textbar{}---\textbar{}---\textbar{}---\textbar{}---\textbar{}\\
\textbar{} R1 (RG) \textbar{} 1.27 k \textbar{} 0.1 \%, 25 ppm/degC \textbar{} sets G1 = 39.90; its TCR adds directly to the AD8226's -100 ppm/degC gain drift \textbar{}\\
\textbar{} R2 / R3 \textbar{} 121 k / 10.0 k \textbar{} 0.1 \%, 25 ppm/degC \textbar{} sets the 0.252 V pedestal; ratio TCR mismatch is a 1.1 uV RTI drift term \textbar{}\\
\textbar{} R4 / R5 \textbar{} 24.9 k / 10.0 k \textbar{} 0.1 \%, 25 ppm/degC \textbar{} sets G2 = 3.49; ratio TCR mismatch adds to gain drift \textbar{}\\
\textbar{} C1, C3, C4 \textbar{} 100 nF \textbar{} X7R, 16 V+ \textbar{} supply and reference bypass \textbar{}\\
\textbar{} C2 \textbar{} 10 uF \textbar{} X5R/X7R, 16 V+ \textbar{} bulk at the power entry (AD8226 Figure 61) \textbar{}
\end{quote}
Absolute tolerance is not what buys accuracy here (zero and span are calibrated downstream per Q7) - TCR is. If P3 finds 1 \% / 25 ppm parts cheaper and in stock, they are acceptable everywhere; 100 ppm/degC thick film is NOT, because it would triple the pedestal drift term and add \textasciitilde{}50 ppm/degC to the gain drift.

Full architecture narrative (not inlined; see the artifact index): \texttt{architecture/blocks.md}, \texttt{architecture/power\_tree.md}, \texttt{architecture/stackup.md}.
\section{Schematic}
\emph{Pending --- produced at P4.}
\section{Layout}
\emph{Pending --- produced at P6.}
\section{Verification}
\emph{Pending --- produced at P8.}
\section{DFM and Fabrication}
\emph{Pending --- produced at P9.}
\section{Run Record (Appendix)}
\subsection*{Phase digests}
\paragraph*{P0}
\begin{quote}\small\ttfamily\raggedright
\# P0 Intake - digest\\
~\\
- Mode `learning block-basics:` -\textgreater{} scope `block-only`, binding `canonical`\\
  (geometry is an OUTPUT; P5 opens `--outline auto`).\\
- `requirements.md` written; `check\_requirements.py` pass, 0 violations,\\
  sections 1-9, mode named. (opus/xhigh - fable tier unavailable, logged.)\\
- 10 questions batched in one round, all answered, appended as section 9a.\\
- Design-changing answers: excitation 3.3 V shared-ground (CM 1.65 V, so a\\
  single-supply in-amp stays viable); output span \textasciitilde{}0.05-3.25 V with a positive\\
  pedestal (the brief's literal 0-3.3 V loses to its own rail); 12-bit usable\\
  calibrated (\textasciitilde{}5 uV RTI = drift + noise); flat within 1\% at 1 kHz.\\
- Derived and binding: \textasciitilde{}7 kHz -3 dB corner =\textgreater{} \textasciitilde{}1 MHz gain-bandwidth at gain\\
  \textasciitilde{}150. Dominant part constraint; rules out low-power zero-drift in-amps.\\
- Carried to P2: the pedestal needs a low-impedance REF drive, or an in-amp\\
  with a high-impedance REF - a divider on a 3-op-amp REF kills CMRR.\\
- Owner delegated remaining judgment; H1-H4 are records, run stops at P9.\\
- Safety: no flags apply (3.3 V SELV bench board, milliamps).
\end{quote}
\paragraph*{P1}
\begin{quote}\small\ttfamily\raggedright
\# P1 Research - digest\\
~\\
- Roster: component-scout (`inamp`) + reference-design (`bridge-front-end`).\\
  Power-architect skipped (trivially powered); interface-spec skipped (screw\\
  terminals are not standards-bound).\\
- REF-pin impedance is the architecture fork, cross-checked TI+ADI: classic\\
  3-op-amp in-amps (AD8226/INA333/INA828) need REF driven low-Z (AD8226 Fig.59\\
  marks the bare divider "INCORRECT"); ICF parts (AD8237/AD8420) have 100M-800M\\
  REF and tolerate a divider.\\
- Bandwidth vs drift squeeze at 3.3 V: INA333 makes drift (0.1 uV/C) but only\\
  \textasciitilde{}2.2 kHz at G\textasciitilde{}150; INA828 has the bandwidth but needs 4.5 V min. AD8226\\
  (2.2-36 V, \textasciitilde{}13.6 kHz at G=147, CMRR 120 dB min, 22-24 nV/rtHz) clears both.\\
- Discrete 3-op-amp killed quantitatively: TI SBOA582 Eq.7 puts 0.01\% resistors\\
  \textasciitilde{}44 dB short of the \textasciitilde{}118 dB this board needs.\\
- RECONCILED BY ORCHESTRATOR: the scout's grounded-REF split dies on AD8226's\\
  (V-)+0.1 V output floor - at G1=100 the bottom 2 mV of span sits in stage-1\\
  saturation. A positive buffered REF is unavoidable; handed to P2 with the math.\\
- P2 must verify: BW-at-gain (interpolated), real RTI drift (VOSI+VOSO/G), and\\
  the Vout-vs-Vcm diamond plot at Vs=3.3 V.
\end{quote}
\paragraph*{P2}
\begin{quote}\small\ttfamily\raggedright
\# P2 Architecture - digest\\
~\\
- Chain: J1 3-pole -\textgreater{} U1 AD8226 G1=39.9 (RG 1.27k, REF on a buffered 0.252 V\\
  node) -\textgreater{} U2B OPA2333-class G2=3.49 whose gain resistor RETURNS TO /VREF -\textgreater{} J2.\\
  U2A buffers the R2/R3 divider. 14 parts, 2 ICs, one sheet.\\
  Vout = 0.252 + 139.2{*}Vdiff. f-3dB 41 kHz (5.9x spec), 0.65 mA, 2.3 mW.\\
  Stackup JLC2313\_1.6, 2 layers. Cost \textasciitilde{}USD 6.6/board.\\
- The split is forced by the DIAMOND PLOT, not by the REF question, and the\\
  claim was VERIFIED on the page: printed p.19 gives Vout = G{*}Vdiff + Vref, so\\
  Eq.2 becomes Vcm + \textbar{}Vout-Vref\textbar{}/2 \textless{} Vs - V+LIMIT with NO gain term. Fig.9\\
  (G=1) and Fig.12 (G=100) are the same hexagon. Max single-stage gain 66-95,\\
  not the 147 the orchestrator proposed or the 100 the scout proposed.\\
- ACCURACY SPEC MISSED AND RECORDED: 13.9 uV typ / 56.4 uV max offset drift\\
  over +-25 C vs a 5 uV budget, plus 30-87 uV gain drift at full scale. That\\
  budget is 10 ppm/degC of full scale - lab-grade - and no 3.3 V-capable part\\
  reaches it. The INA333 alternative was costed and rejected on evidence.\\
- Mode relaxed nothing (no size is stated anywhere).\\
- RESEARCH: 4 tasks, 2 rounds, 6 of 6 budget spent. Round 1 seeded the domain\\
  (library held only buck records) but was witnessed by the wrong part -\\
  7 of 13 records refuted. Round 2 re-witnessed from the AD8226 itself via an\\
  allowlisted Farnell mirror: 9 of 9 verified. Net 16 verified records + 2 new\\
  checklists (`inamp`, `in`).\\
- P2 exit coverage: exit 0, 2 slots, 0 gaps, all 11 classes provisional.\\
  Provisional is the correct ceiling - only the owner's promotion makes them\\
  approved, only bring-up makes them proven.\\
- 8 sim benches specified with numeric windows for the P8 `sim` gate.
\end{quote}
\subsection*{Gate attempt history}
\emph{(no entries)}
\subsection*{Issues}
\emph{(no entries)}
\subsection*{Human checkpoints}
{\small
\begin{longtable}{lllp{2.6cm}p{6.4cm}}
\toprule
\textbf{Id} & \textbf{Phase} & \textbf{Status} & \textbf{When} & \textbf{Note} \\
\midrule
\endhead
1 & P2 & not reached &  &  \\
2 & P4 & not reached &  &  \\
3 & P6 & not reached &  &  \\
4 & P8 & not reached &  &  \\
5 & P10 & not reached &  &  \\
\bottomrule
\end{longtable}
}
\section{Artifact Index}
{\small
\begin{longtable}{p{9cm}rp{4cm}}
\toprule
\textbf{File} & \textbf{Size} & \textbf{Note} \\
\midrule
\endhead
\texttt{state.json} & 45,557 B &  \\
\texttt{brief/brief.md} & 575 B &  \\
\texttt{requirements.md} & 17,532 B &  \\
\texttt{architecture/blocks.md} & 27,373 B &  \\
\texttt{architecture/constraints.json} & 3,691 B &  \\
\texttt{architecture/power\_tree.md} & 4,310 B &  \\
\texttt{architecture/sheets.md} & 4,394 B &  \\
\texttt{architecture/stackup.md} & 3,813 B &  \\
\bottomrule
\end{longtable}
}
\end{document}
